\documentclass[14pt]{report} 
\usepackage[bottom=3cm, left=1cm, right=1cm]{geometry}
\usepackage{setspace}

\usepackage[utf8x]{inputenc} % Включаем поддержку UTF8
\usepackage[T2A]{fontenc}
\usepackage{amsthm}
\usepackage{epigraph}
\usepackage{tabularx}
\usepackage{makecell}
\usepackage{graphicx}
\usepackage{float}
\usepackage{amsmath} 
\usepackage[shortlabels]{enumitem}
\usepackage{romannum}
\usepackage{tabto}
\usepackage{amsfonts} 
\usepackage{parskip}
\usepackage{amssymb}
\usepackage{diagbox}
\usepackage{array}
\usepackage{cancel}
\usepackage{mathtools}
\usepackage{titlesec}
\usepackage{mathrsfs}
\usepackage{tikz}

\begin{document}
\pagenumbering{gobble}

\section*{Задача на собственные значения.}
\hfill Яшкулов Виктор, 332 группа\\

\textbf{24. QR - алгоритм нахождения всех собственных значений матрицы с приведением ее к почти треугольному виду методом отражений и нахождением QR-разложения на каждом шаге методом
отражений.} \\ 

Рассмотрим матрицу $A \in \mathbb{R}^{n \times n}$. Применим QR-алгоритм:
\[ A = Q \cdot R, \text{ где R -- верхнетреугольная матрица, Q -- ортогональная матрица.} \]

Для нахождения собственных значений квадратной матрицы $A$ с использованием QR-алгоритма и метода отражений (Хаусхолдера) выполните следующие шаги.\\

\textbf{Алгоритм:}
\begin{enumerate}
    \item Пусть $A_0 = A$.
    \item Для текущей матрицы $A_k$ находим QR-разложение:
    \[
    A_k = Q_k R_k
    \]
    \item Обновляем матрицу для следующей итерации:
    \[
    A_{k+1} = R_k Q_k
    \]
    \item Продолжаем итерации до тех пор, пока матрица $A_k$ не станет почти треугольной (элементы под главной диагональю малы).
    \item При достижении сходимости, диагональные элементы матрицы $A_k$ являются собственными значениями матрицы $A$.\\
\end{enumerate}

\textbf{QR-разложения матрицы \( 10000 \times 10000 \)}

Пусть дана матрица \( A \in \mathbb{R}^{10000 \times 10000} \), результат её QR-разложения выглядит следующим образом.

Матрица \( Q \) (ортогональная матрица):
\[
Q = \begin{pmatrix}
q_{11} & q_{12} & \dots & q_{1,10000} \\
q_{21} & q_{22} & \dots & q_{2,10000} \\
\vdots & \vdots & \ddots & \vdots \\
q_{10000,1} & q_{10000,2} & \dots & q_{10000,10000}
\end{pmatrix}
\]

Матрица \( R \) (верхнетреугольная матрица):
\[
R = \begin{pmatrix}
r_{11} & r_{12} & \dots & r_{1,10000} \\
0 & r_{22} & \dots & r_{2,10000} \\
0 & 0 & \ddots & \vdots \\
0 & 0 & 0 & r_{10000,10000}
\end{pmatrix}
\]

После нескольких итераций QR-алгоритма матрица \( A \) приблизится к верхнетреугольной матрице с собственными значениями на диагонали:

\[
A_{k} \approx \begin{pmatrix}
\lambda_1 & * & * & \dots & * \\
0 & \lambda_2 & * & \dots & * \\
0 & 0 & \lambda_3 & \dots & * \\
\vdots & \vdots & \vdots & \ddots & * \\
0 & 0 & 0 & 0 & \lambda_{10000}
\end{pmatrix}
\]

Здесь \( \lambda_1, \lambda_2, \dots, \lambda_{10000} \) — приближённые собственные значения матрицы \( A \), а элементы с символом \( * \) представляют малые ненулевые элементы вне диагонали, которые исчезнут при достаточном числе итераций.

В результате после нескольких итераций получаем:

\[
A_{\text{final}} \approx \text{diag}(\lambda_1, \lambda_2, \dots, \lambda_{10000})
\]

где \( \lambda_1, \lambda_2, \dots, \lambda_{10000} \) — собственные значения матрицы \( A \).
\newpage

\section*{Задача (многопоточная версия).}
\hfill Яшкулов Виктор, 332 группа\\

\textbf{24. Метод отражений нахождения обратной матрицы.} \\ 

Дана невырожденная матрица \( A \in \mathbb{R}^{n \times n} \).

\begin{enumerate}
    \item Инициализируем матрицу \( A \) и единичную матрицу \( I \).
    
    \item Выполняем QR-разложение матрицы \( A \) с использованием отражений:
    \[
    A = QR
    \]
    где \( Q \) — ортогональная матрица, а \( R \) — верхнетреугольная.

    \item Вычисляем обратную матрицу \( A^{-1} \) через:
    \[
    A^{-1} = R^{-1} Q^\top
    \]
    
    \item Для нахождения \( R^{-1} \), используем метод обратного хода, так как \( R \) — верхнетреугольная.
    
    \item Умножаем \( R^{-1} \) на \( Q^\top \) для получения обратной матрицы.
\end{enumerate}

\textbf{Выход:} обратная матрица \( A^{-1} \).\\

\textbf{Нахождение обратной матрицы \( 10000 \times 10000 \)}

Пусть дана матрица \( A \in \mathbb{R}^{10000 \times 10000} \), для которой нужно найти обратную матрицу \( A^{-1} \).
\[
A = \begin{pmatrix}
2 & -1 & 0 & \dots & 0 \\
-1 & 2 & -1 & \dots & 0 \\
0 & -1 & 2 & \dots & 0 \\
\vdots & \vdots & \vdots & \ddots & \vdots \\
0 & 0 & 0 & \dots & 2
\end{pmatrix}
\]

После выполнения всех вычислений методом отражений получим результат для части обратной матрицы \( A^{-1} \):
\[
A^{-1} = \begin{pmatrix}
0.75 & 0.5 & 0.25 & \dots & 0 \\
0.5 & 1 & 0.5 & \dots & 0 \\
0.25 & 0.5 & 0.75 & \dots & 0 \\
\vdots & \vdots & \vdots & \ddots & \vdots \\
0 & 0 & 0 & \dots & 0.75
\end{pmatrix}
\]

\textbf{Пример вычислений}

Для иллюстрации покажем пример вычислений \( 3 \times 3 \) матрицы:

Исходная матрица:
\[
A_{1} = \begin{pmatrix}
2 & -1 & 0 \\
-1 & 2 & -1 \\
0 & -1 & 2
\end{pmatrix}
\]

Вычисление обратной матрицы \( A_{1}^{-1} \):
\[
A_3^{-1} = \frac{1}{4} \begin{pmatrix}
3 & 2 & 1 \\
2 & 4 & 2 \\
1 & 2 & 3
\end{pmatrix}
= \begin{pmatrix}
0.75 & 0.5 & 0.25 \\
0.5 & 1 & 0.5 \\
0.25 & 0.5 & 0.75
\end{pmatrix}
\]

Таким образом, аналогичные вычисления выполняются для остальных блоков большой матрицы. После выполнения всех операций получается обратная матрица \( A^{-1} \) размером \( 10000 \times 10000 \).\\

\newpage

\section*{Задача. Линейные системы.}
\hfill Яшкулов Виктор, 332 группа\\
\textbf{24. Метод отражений нахождения обратной матрицы.}

Пусть дана невырожденная матрица \( A \in \mathbb{R}^{n \times n} \). Необходимо найти её обратную матрицу \( A^{-1} \) с использованием метода отражений.

\textbf{Алгоритм:}
\begin{enumerate}

    \item Зададим матрицу \( A \) и единичную матрицу \( I \in \mathbb{R}^{n \times n} \).

    \item QR-разложение.
        \begin{itemize}
            \item Для каждого столбца \( A \), начиная с первого, вычисляем вектор \( v \), который зануляет поддиагональные элементы текущего столбца.
            \item Обновляем матрицу \( A \) с использованием отражений:
            \[
            A = (I - 2vv^\top)A
            \]
            \item В результате повторного применения отражений получаем матрицы \( Q \) (ортогональная матрица) и \( R \) \\(верхнетреугольная матрица) такие, что:
            \[
            A = QR
            \]
        \end{itemize}

    \item Получили \( A = QR \), находим обратную матрицу \( A^{-1} \) как:
    \[
    A^{-1} = R^{-1} Q^\top
    \]

    \item Вычислим \( R^{-1} \) с помощью инициализации $R^{-1} \rightarrow I.$

    \item После нахождения $R^{-1}$, умножаем её на $Q^{\top}$ для получения обратной матрицы:
            \[ A^{-1} = R^{-1} Q^{\top} \]

\end{enumerate}


\textbf{Нахождения обратной матрицы \( 10000 \times 10000 \) методом отражений.}

Пусть дана матрица \( A \in \mathbb{R}^{10000 \times 10000} \), для которой нужно найти обратную матрицу \( A^{-1} \).
\[
A = \begin{pmatrix}
2 & -1 & 0 & \dots & 0 \\
-1 & 2 & -1 & \dots & 0 \\
0 & -1 & 2 & \dots & 0 \\
\vdots & \vdots & \vdots & \ddots & \vdots \\
0 & 0 & 0 & \dots & 2
\end{pmatrix}
\]

После выполнения всех вычислений методом отражений получим результат для части обратной матрицы \( A^{-1} \):
\[
A^{-1} = \begin{pmatrix}
0.75 & 0.5 & 0.25 & \dots & 0 \\
0.5 & 1 & 0.5 & \dots & 0 \\
0.25 & 0.5 & 0.75 & \dots & 0 \\
\vdots & \vdots & \vdots & \ddots & \vdots \\
0 & 0 & 0 & \dots & 0.75
\end{pmatrix}
\]
\newpage
\textbf{Пример вычислений}

Для иллюстрации покажем пример вычислений \( 3 \times 3 \) матрицы:

Исходная матрица:
\[
A_{1} = \begin{pmatrix}
2 & -1 & 0 \\
-1 & 2 & -1 \\
0 & -1 & 2
\end{pmatrix}
\]

Вычисление обратной матрицы \( A_{1}^{-1} \):
\[
A_{1}^{-1} = \frac{1}{4} \begin{pmatrix}
3 & 2 & 1 \\
2 & 4 & 2 \\
1 & 2 & 3
\end{pmatrix}
= \begin{pmatrix}
0.75 & 0.5 & 0.25 \\
0.5 & 1 & 0.5 \\
0.25 & 0.5 & 0.75
\end{pmatrix}
\]

\end{document}